\documentclass[aps,prd,10pt,superscriptaddress,nofootinbib,showkeys,amsmath,amssymb]{revtex4-2}

\usepackage{graphicx}
\usepackage{bm}
\usepackage{booktabs}
\usepackage{array}
\usepackage{amsthm}
\usepackage{enumerate}
\usepackage{hyperref}
\usepackage{lineno}

\hypersetup{
  colorlinks=true,
  linkcolor=blue,
  citecolor=blue,
  urlcolor=blue
}

\newtheorem{hypothesis}{Hypothesis}
\newtheorem{prediction}{Prediction}
\newtheorem{criterion}{Criterion}
\newcolumntype{L}[1]{>{\raggedright\arraybackslash}p{#1}}
\setlength{\emergencystretch}{2em}
\newcommand{\flux}{\Phi}
\newcommand{\constraint}{C}
\newcommand{\phase}{\Psi_s}

\begin{document}

\title{Paper XII: Friction, Superconductivity, Plasma States, and Transport Mysteries: A Constraint-Driven Flux Map for Collective Matter}

\author{Steve Bico Mujjabi, MD}
\email{msbico@gmail.com}
\homepage{https://github.com/bampita-bico/CDFD-Part-I-Release}
\thanks{Independent researcher; ORCID: \href{https://orcid.org/0009-0001-0556-5516}{0009-0001-0556-5516}.}
\affiliation{VuraLabs, Kampala, Uganda}

\date[]{July 2026}

\begin{abstract}
Paper XI extended CDFD into cosmological mysteries. Paper XII turns the same discipline toward laboratory-scale and engineering-scale mysteries: friction, stick-slip rupture, superconductivity, superfluidity, plasma confinement, magnetic reconnection, turbulence, jamming, glasses, and anomalous transport. These systems are not treated as solved by analogy. They are treated as falsifiable test beds for one claim: many physical mysteries arise where flux is forced through adaptive constraints whose capacity, memory, and topology change during transport. The paper separates established physics from CDFD interpretation, gives domain-specific mappings for $\Phi$, $C$, and $\Psi_s$, states failure conditions for each proposed bridge, and includes reproducible outputs generated by the Paper XII supplementary script.

\medskip
\noindent In the extended notation used across the series, the vacuum limit is written
as $\Psi_s = (\Phi/C)\cdot S\cdot M_s$. In this manuscript the same notation is
used cautiously as a transport map for measurable collective systems. The public
CDFD Runtime implements this state grammar for reusable simulations;
the paper-local supplementary script remains the audited reproduction layer for
the figures and tables in this Part I archive.
\end{abstract}

\keywords{Constraint-Driven Field Theory, Constraint-Driven Flux Dynamics, adaptive operating ratio, vacuum structure, Mujjabi tests, reproducible physics, friction, superconductivity, plasma states, anomalous transport}

\maketitle
\linenumbers

\section*{Project Context}
This manuscript constitutes Part I (Paper 12) of the multi-part series
\textit{Constraint-Driven Flux Dynamics (CDFD)}. The underlying mathematical
architecture builds directly upon the concepts established in Paper 11 of this
series (Steve Bico Mujjabi, \textit{``CDFD Part I: Fundamental Physics ---
Paper 11: Dark Matter, Dark Energy, Black Holes, and the Blancken-Layer
Synthesis: A Constraint-Driven Flux Interpretation of Cosmological
Anomalies''}, manuscript; paper-specific DOI pending).
Its local reference chain includes the preceding Part I manuscripts, Papers
01--11~\cite{MujjabiI2026,MujjabiII2026,MujjabiIII2026,MujjabiIV2026,
MujjabiV2026,MujjabiVI2026,MujjabiVII2026,MujjabiVIII2026,MujjabiIX2026,
MujjabiX2026,MujjabiXI2026}.

\section*{Symbol Conventions}

\noindent\textbf{CDFD physical-vacuum symbol convention.}
\begin{center}
{\small
\begin{tabular}{L{0.16\linewidth}L{0.33\linewidth}L{0.41\linewidth}}
\hline
Symbol & Part I meaning & Physical reading \\
\hline
$J$ or $\Phi$ & Driving flux or throughput & Transport current, field throughput, circulation drive, or generic flux; a subscript fixes the exact channel \\
$C$ & Active constraint or capacity burden & Vacuum transport capacity, geometric confinement, stiffness, boundary resistance, or load-bearing limit \\
$S$ & Surface responsiveness & Local ability of the vacuum or modeled surface to reconfigure under drive \\
$M_s$ & Structural memory & Retained geometry, phase history, stiffness history, or accumulated constraint state \\
$\Psi_s$ & Adaptive operating ratio $(J/C)\cdot S\cdot M_s$ or $(\Phi/C)\cdot S\cdot M_s$ & Local bookkeeping gauge for constrained operation; it is not a new universal constant \\
$R,a,\chi$ & Ring radius, core radius, and aspect ratio $R/a$ & Geometry used to connect vortex structure to dimensionless electromagnetic scales \\
$\rho_0,c,\kappa$ & Vacuum density scale, propagation speed, and transport stiffness & Medium parameters used in stability, pressure, and self-consistency calculations \\
\hline
\end{tabular}
}
\end{center}

\noindent\textbf{Greek-symbol convention.}
\begin{center}
{\small
\begin{tabular}{L{0.16\linewidth}L{0.33\linewidth}L{0.41\linewidth}}
\hline
Symbol & Local role & Interpretation rule \\
\hline
$\alpha$ & Fine-structure constant unless a local equation states otherwise & In the Part I physics papers, un-subscripted $\alpha$ means the electromagnetic fine-structure constant \\
$\beta,\gamma$ & Local response, stiffness, coupling, or correction coefficients & Equation-local coefficients; subscripts bind them to the named mode, channel, or approximation \\
$\theta,\phi,\Phi_c$ & Phase, orientation, or chirality angles & Used for torus-knot orientation, vacuum phase, or chiral phase geometry as defined in the nearby equation \\
$\delta,\epsilon$ & Perturbation, residual, tolerance, or small-offset scales & Local stability margins, numerical residuals, or experimental tolerances \\
$\lambda,\mu,\nu$ & Rate, mass, mode, or frequency labels & Meaning is fixed by the equation, subscript, and physics context where the symbol appears \\
$\rho,\sigma,\tau$ & Density, spread/noise, or time-scale terms & Local density, variance, stress, relaxation, or transport-memory quantities \\
$\Omega,\Lambda$ & State/domain space or integrated measure where defined & Reserved for explicitly defined state spaces, spectra, or synthesis-level quantities \\
\hline
\end{tabular}
}
\end{center}

\section{Purpose}

The physics series now has two layers. Papers I--X contain the formal CDFD
mass-spectrum and vacuum equation-of-state ladder. Paper XI is the speculative
cosmology bridge. Paper XII adds a more testable bridge: condensed matter and
collective media. Friction, superconductivity, plasma states, turbulence, and
jamming are useful because they are measurable. They can be probed in the
laboratory, simulated, and compared against mature standard theories of contact
mechanics, superconductivity, vortex matter, transition to turbulence, and
granular media~\cite{BowdenTabor1950,BCS1957,Abrikosov1957,Reynolds1883,JaegerNagelBehringer1996}.
This is also the scale where emergence is expected: collective systems can have
organizing variables that are not visible in the microscopic description alone
~\cite{Anderson1972}.

The manuscript therefore uses a strict posture:
\begin{enumerate}
    \item conventional descriptions remain primary;
    \item CDFD is introduced only as an organizing transport map;
    \item every proposed bridge must name a measurable discriminant;
    \item if the discriminant is not found after controlling for standard variables, the bridge fails.
\end{enumerate}

\section{Common Transport Form}

Let a physical system be described by a driving flux $\Phi_X$, a transport constraint $C_X$, and a local operating ratio
\begin{equation}
    \Psi_{s,X}=\frac{\Phi_X}{C_X}.
\end{equation}
For laboratory matter, $C_X$ should not be read as a single universal substance. It is a domain-specific effective constraint: frictional contact strength, dephasing burden, disorder, magnetic confinement capacity, turbulent impedance, or jamming resistance.

The common adaptive law is
\begin{equation}
    \partial_t C_X
    =
    \alpha_X|\Phi_X|
    +\eta_X M_X
    -\beta_X C_X
    +\gamma_X\nabla^2 C_X,
    \label{eq:adaptive_constraint}
\end{equation}
where $M_X$ is a memory, damage, pinning, or ordering field. The added $M_X$ term matters in this paper because frictional aging, vortex pinning, glass relaxation, and plasma hysteresis cannot be represented by a memoryless ratio.
The coefficients are domain-specific: $\alpha_X$ reinforces the effective
constraint in response to flux, $\eta_X$ weights the memory or ordering field,
$\beta_X$ relaxes the constraint, and $\gamma_X$ smooths or diffuses the
constraint field in space.

\subsection{Computational Status of This Paper}

The figures in this manuscript are not empirical fits. They are deterministic toy outputs whose purpose is to make the proposed mappings explicit enough to fail. The generator \texttt{supplementary\_XII.py} writes:
\begin{itemize}
    \item a transport-mystery map linking each domain to its flux, constraint, transition signal, and failure condition;
    \item a stick-slip memory model with detected release events;
    \item a superconducting constraint-reallocation curve;
    \item a plasma rupture-index field;
    \item a jamming memory distribution;
    \item symbolic and numerical gates checking that the examples behave as stated.
\end{itemize}
The outputs therefore serve as a scaffold for experiments, not as evidence that the bridge is already true.

\begin{figure}[ht]
\centering
\includegraphics[width=0.98\textwidth]{figures/fig12_0_transport_mystery_map.pdf}
\caption{Generated Paper XII mystery map. The CDFD claim is not that one scalar replaces existing theories, but that many transitions are organized by adaptive transport constraints. Source table: \texttt{table1\_transport\_mystery\_map.csv}.}
\label{fig:transport_mystery_map}
\end{figure}

\section{Friction and Stick-Slip}

Friction is often introduced as a force opposing motion, but real friction is not merely a number. It has memory, roughness, contact aging, heating, wear, acoustic emission, and abrupt slip. This makes it a natural CDFD target.

Let the shear flux be
\begin{equation}
    \Phi_f = \tau v A,
\end{equation}
where $\tau$ is shear stress, $v$ is slip velocity, and $A$ is the apparent contact area. Let the frictional constraint be
\begin{equation}
    C_f = \mu_{\mathrm{eff}} N A_r,
\end{equation}
where $N$ is normal load, $A_r$ is real contact area, and $\mu_{\mathrm{eff}}$ includes asperity chemistry, roughness, lubrication, temperature, and wear state.

Stick-slip occurs when the system crosses a constraint boundary:
\begin{equation}
    \Psi_f=\frac{\tau v A}{\mu_{\mathrm{static}}N A_r}
    \rightarrow 1^{+}.
\end{equation}
During sticking, $A_r$ and memory $M_f$ can grow:
\begin{equation}
    \partial_t M_f = a_f N - b_f v M_f - h_f T M_f.
\end{equation}
During slip, the constraint partially collapses:
\begin{equation}
    C_f^{\mathrm{dynamic}} < C_f^{\mathrm{static}}.
\end{equation}

\begin{hypothesis}[Friction as adaptive constraint memory]
Frictional behavior is governed not only by instantaneous load and velocity, but by an adaptive constraint field whose memory variable controls the transition between static locking, dynamic sliding, and wear-mediated reorganization.
\end{hypothesis}

\begin{prediction}[Frictional precursor]
Before macroscopic slip, acoustic emission, microscopic rearrangement, or contact-resistance noise should increase preferentially in regions where $\nabla(C_f^{-1}\nabla\Phi_f)$ is large, not merely where absolute load is maximal.
\end{prediction}

\begin{criterion}[Friction falsification]
The CDFD friction bridge fails if stick-slip timing and localization are fully explained by standard rate-and-state or contact-mechanics variables with no residual dependence on a transport-memory field.
\end{criterion}

\begin{figure}[ht]
\centering
\includegraphics[width=0.92\textwidth]{figures/fig12_1_friction_stick_slip.pdf}
\caption{Generated stick-slip output. A slowly increasing shear-flux proxy loads an adaptive contact constraint until release. The lower panel shows a precursor/noise proxy. Event table: \texttt{table2\_friction\_stick\_slip\_events.csv}.}
\label{fig:friction_stick_slip}
\end{figure}

\subsection{Why Friction Matters for CDFD}

Friction is strategically important because it is the simplest everyday case where constraint is visibly real. A surface can be immobile under stress, then fail abruptly, then rebuild its contact structure. The same material pair may behave differently after waiting, heating, polishing, contaminating, or repeatedly sliding. That history dependence is exactly the kind of memory term represented by $M_X$ in Eq.~\eqref{eq:adaptive_constraint}.

The CDFD addition to standard tribology is not a new coefficient of friction. It is a localization claim: slip should nucleate where flux, constraint weakness, and constraint gradients combine, not necessarily where load alone is largest. This gives the bridge a practical experimental target: synchronized contact imaging, acoustic emission, and local shear maps should identify whether $\nabla(C_f^{-1}\nabla \Phi_f)$ adds predictive value.

\subsection{Triboelectricity}

Triboelectric charging belongs in the same family. Contact and separation do not merely rub charges loose; they create and destroy interfacial transport channels. A CDFD formulation would track a charge flux $\Phi_q$ and an interfacial constraint $C_q$ set by surface chemistry, humidity, roughness, trapped states, and separation rate. The useful claim is narrow: charge sign, magnitude, and relaxation should show path dependence on prior constraint history. If charging is completely predicted from instantaneous material pairs and environment, the CDFD memory bridge adds nothing.

\section{Superconductivity and Superfluidity}

Superconductivity is the cleanest place to avoid sloppy language. If $C$ is called ``constraint'', one might incorrectly say superconductivity is $C\to0$. That is not right. A superconductor has vanishing dissipative resistance, but it also has strong phase rigidity, quantized flux, coherence length, penetration depth, vortex pinning, and critical current.

Paper XII therefore separates two constraints:
\begin{equation}
    C_{\mathrm{diss}} =
    C_{\mathrm{phonon}}+C_{\mathrm{disorder}}+C_{\mathrm{dephase}}+C_{\mathrm{vortex}},
\end{equation}
and
\begin{equation}
    C_{\mathrm{phase}} = \rho_s,
\end{equation}
where $\rho_s$ is the superfluid or superconducting phase stiffness. The superconducting transport index is then
\begin{equation}
    \mathcal{S}_{\mathrm{SC}}
    =
    \frac{\Phi_{\mathrm{pair}}\,C_{\mathrm{phase}}}
         {C_{\mathrm{diss}}+\epsilon}.
\end{equation}
The superconducting transition is a regime change in which dissipative constraint collapses while phase constraint rises.

\begin{figure}[ht]
\centering
\includegraphics[width=0.92\textwidth]{figures/fig12_2_superconductivity_reallocation.pdf}
\caption{Generated superconductivity output. The toy model separates dissipative constraint from phase stiffness. Below the transition, dissipative resistance falls while phase/topological constraint remains active. Curve table: \texttt{table3\_superconductivity\_reallocation\_curve.csv}.}
\label{fig:superconductivity_reallocation}
\end{figure}

\begin{hypothesis}[Coherence as constraint reallocation]
Superconductivity is a constraint reallocation transition: dissipative scattering channels are suppressed, while phase stiffness and topological flux constraints become the dominant regulators of transport.
\end{hypothesis}

This framing is compatible with conventional superconductivity only if it reproduces known limits: ordinary BCS superconductors, vortex dynamics in type-II materials, Josephson phase coherence, and critical current. It becomes interesting only where ordinary explanation is incomplete, especially high-$T_c$ superconductors, strange metals, pseudogap regimes, and disorder-driven superconductor-insulator transitions.

\subsection{Vortices and Pinning}

In type-II superconductors, dissipation often appears through vortex motion. The relevant constraint is not ordinary resistance alone, but pinning topology:
\begin{equation}
    C_{\mathrm{vortex}} =
    C_{\mathrm{pin}}(\mathbf{x})
    - C_{\mathrm{drive}}(J,B,T).
\end{equation}
When current drive exceeds pinning capacity, vortices move and resistance returns. This gives a falsifiable CDFD statement: vortex avalanches and flux jumps should localize where the transport operator
\begin{equation}
    \nabla\cdot(C_{\mathrm{vortex}}^{-1}\nabla\Phi_J)
\end{equation}
is extremal.

\begin{prediction}[Superconducting discriminant]
Across a superconductor-insulator transition, collapse of resistance should be accompanied by a measurable reallocation from dissipative impedance to phase stiffness or pinning topology. If resistance falls without any measurable change in coherence, stiffness, or topological constraint, the CDFD interpretation is weakened.
\end{prediction}

\begin{criterion}[Superconductivity falsification]
The superconductivity bridge fails if it cannot recover standard superconducting constraints and if its added variables do not improve prediction of high-$T_c$, vortex, or disorder-driven anomalies.
\end{criterion}

\subsection{Superfluidity}

Superfluidity is the neutral analogue: mass flow proceeds without ordinary viscous dissipation, but quantized vortices and phase coherence still constrain transport. The CDFD reading is again not ``no constraint''. It is a switch from dissipative constraint to topological phase constraint.

\subsection{\texorpdfstring{High-$T_c$ and Strange-Metal Boundary}{High-Tc and Strange-Metal Boundary}}

High-$T_c$ superconductivity is where Paper XII could eventually matter most. Conventional superconductivity already has a mature theory. The difficult territory is where pairing, phase stiffness, disorder, stripe order, spin fluctuations, pseudogap behavior, and strange-metal transport do not collapse into one simple transition.

In CDFD language, this is a decoupling problem:
\begin{equation}
    C_{\mathrm{pair}}\neq C_{\mathrm{phase}}\neq C_{\mathrm{diss}}.
\end{equation}
Pair formation may begin while phase coherence remains constrained; dissipation may remain anomalous even when a partial gap exists. A useful CDFD model must therefore predict which constraint collapses first and which remains active. If it cannot distinguish pairing onset from phase locking, it adds no value.

\section{Plasma States and Magnetic Reconnection}

Plasma systems already fit CDFD language naturally. A plasma is a medium where charged-particle flux, magnetic topology, pressure gradients, turbulence, and radiation losses compete. In a confinement device, solar corona, or magnetosphere, the relevant variables are:
\begin{align}
    \Phi_p &\sim \text{current density, pressure-gradient drive, or energy flux},\\
    C_p &\sim \text{magnetic topology, confinement capacity, resistivity, and turbulent transport}.
\end{align}

The public model used here states only a scientific rupture criterion: plasma
instability is expected when a high-$\Psi$ drive coincides with steep constraint
gradients. In this manuscript the criterion is written as
\begin{equation}
    \mathcal{R}_p =
    \left|\nabla C_p\right|
    \frac{\Phi_p}{C_p}
\end{equation}
is a candidate rupture index for reconnection, edge-localized modes, or disruption-like events.

\begin{figure}[ht]
\centering
\includegraphics[width=0.98\textwidth]{figures/fig12_3_plasma_rupture_index.pdf}
\caption{Generated plasma output. The rupture index combines drive, constraint, and constraint-gradient information. In the toy field, risk does not follow the drive field alone. Summary table: \texttt{table5\_plasma\_rupture\_index\_summary.csv}.}
\label{fig:plasma_rupture}
\end{figure}

\begin{hypothesis}[Plasma rupture as constraint-topology failure]
Plasma disruptions and reconnection events occur where current or pressure-gradient flux exceeds the local capacity of magnetic topology to preserve confinement.
\end{hypothesis}

\begin{prediction}[Plasma precursor]
Before a disruption, the highest-risk region should be better predicted by a combined drive-plus-topology index such as $\mathcal{R}_p$ than by absolute plasma energy density alone.
\end{prediction}

\begin{criterion}[Plasma falsification]
The plasma bridge fails if disruptions, reconnection, or confinement loss correlate only with standard scalar plasma parameters and not with constrained transport geometry after controlling for known MHD variables.
\end{criterion}

\subsection{Plasma States as CDFD Regimes}

\begin{center}
\begin{tabular}{L{0.22\textwidth}L{0.28\textwidth}L{0.34\textwidth}}
\toprule
\textbf{Regime} & \textbf{CDFD reading} & \textbf{Observable handle}\\
\midrule
Diffuse plasma & low coherence, weak constraint topology & broad particle/field distribution\\
Confined plasma & $\Psi_p\approx1$ with magnetic topology sustaining flux & long confinement time, low loss rate\\
Reconnection & steep constraint gradient breaks and reroutes flux & rapid magnetic-energy release\\
Runaway/\allowbreak disruption & sustained $\Psi_p>1$ over critical region & quench, edge loss, global instability\\
Ignition-like state & self-heating flux maintains organized constraint & sustained fusion-relevant balance\\
\bottomrule
\end{tabular}
\end{center}

\section{Turbulence, Drag, and Boundary Layers}

Turbulence is the classical transport mystery. A laminar flow has organized flux and stable constraint. A turbulent flow reorganizes energy across scales. CDFD should not pretend to replace Navier--Stokes. Its role is to provide a phase-structure language:
\begin{equation}
    \Psi_{\mathrm{fluid}}
    =
    \frac{\Phi_{\mathrm{inertial}}}
         {C_{\mathrm{viscous}}+C_{\mathrm{boundary}}+C_{\mathrm{eddy}}}.
\end{equation}
The denominator is intentionally adaptive: once turbulence begins, eddy structure becomes a new transport constraint rather than a passive consequence.

\begin{hypothesis}[Turbulence as adaptive constraint production]
Turbulence is a regime in which inertial flux produces its own multiscale constraint field, redistributing transport through eddies, boundary layers, and coherent structures.
\end{hypothesis}

\begin{prediction}[Transition localization]
The laminar-to-turbulent transition should localize where boundary-layer constraint cannot dissipate inertial flux smoothly, not simply where global Reynolds number is high.
\end{prediction}

\begin{criterion}[Turbulence falsification]
The turbulence bridge fails if no local constraint-memory or geometry variable improves transition prediction beyond standard hydrodynamic variables.
\end{criterion}

\subsection{Drag Reduction and Flow Control}

Drag reduction is a practical reason to keep turbulence in the mystery queue. Polymers, riblets, compliant walls, bubbles, and active flow-control systems all alter the way flux encounters boundary constraint. CDFD predicts that successful drag reduction should not merely lower average resistance; it should change where constraint gradients form and how coherent structures reroute momentum.

This is experimentally concrete. Particle-image velocimetry and wall-shear sensors can compare three fields: inertial flux, boundary constraint, and coherent-structure transport. If drag reduction works while those fields remain unchanged except for a scalar viscosity shift, the CDFD interpretation is weak. If it works by redistributing constraint topology, the bridge becomes more interesting.

\section{Glasses, Jamming, Granular Flow, and Avalanches}

Glasses and jammed materials are almost the inverse of superconductors. Superconductors suppress dissipative constraint and permit coherent transport. Jammed systems accumulate constraint faster than flux can relax it.

For a granular or glassy material:
\begin{equation}
    \Phi_g \sim \text{imposed shear, vibration, or thermal agitation},
\end{equation}
and
\begin{equation}
    C_g \sim \text{cage topology, packing fraction, contact network, and relaxation barrier}.
\end{equation}
The jamming boundary is
\begin{equation}
    \Psi_g=\frac{\Phi_g}{C_g}\ll1
\end{equation}
for trapped, over-constrained states, while avalanche failure occurs when stored stress abruptly exceeds local constraint:
\begin{equation}
    \Psi_g^{\mathrm{stored}}\rightarrow1^{+}.
\end{equation}

\begin{figure}[ht]
\centering
\includegraphics[width=0.92\textwidth]{figures/fig12_4_jamming_memory.pdf}
\caption{Generated jamming output. Two systems with the same nominal drive but different preparation memory produce different avalanche tails. Distribution table: \texttt{table6\_jamming\_memory\_distributions.csv}.}
\label{fig:jamming_memory}
\end{figure}

\begin{prediction}[Jamming memory]
Two samples with the same density and shear rate but different preparation histories should show different avalanche statistics if the CDFD memory field $M_g$ is real.
\end{prediction}

\begin{criterion}[Jamming falsification]
The jamming bridge fails if preparation history, contact-network memory, and constraint topology add no predictive information beyond instantaneous density, stress, and temperature.
\end{criterion}

\section{Fracture, Cracks, and Material Failure}

Fracture is friction's irreversible cousin. A crack is a transport event: elastic energy flux localizes into a narrowing process zone while material constraint weakens at the crack tip. In a CDFD representation:
\begin{align}
    \Phi_{\mathrm{fracture}} &\sim \text{elastic energy release rate},\\
    C_{\mathrm{fracture}} &\sim \text{local toughness, plastic dissipation, microstructural barriers}.
\end{align}
The crack path should follow corridors where stored energy flux overwhelms local material constraint:
\begin{equation}
    \mathcal{K}_{\mathrm{CDFD}}
    =
    \nabla\cdot\left(C_{\mathrm{fracture}}^{-1}\nabla \Phi_{\mathrm{fracture}}\right).
\end{equation}
This does not replace stress-intensity factors or fracture mechanics. It asks whether heterogeneous fracture paths, branching, and avalanche acoustic emission are better predicted when constraint topology is explicitly mapped.

\begin{prediction}[Fracture corridor]
In heterogeneous materials, crack branching should align with high $\mathcal{K}_{\mathrm{CDFD}}$ corridors after controlling for ordinary stress intensity.
\end{prediction}

\begin{criterion}[Fracture falsification]
The fracture bridge fails if crack paths and acoustic-emission bursts are fully predicted from standard fracture variables with no residual role for adaptive constraint topology.
\end{criterion}

\section{Topological and Field-Written Matter}

Several modern materials are controlled by topology or externally written constraint fields. Quantum Hall systems, topological insulators, ferrofluids, magnetorheological fluids, and vortex matter in superconductors all show transport that is not just local resistance. Boundaries, defects, and applied fields determine which channels are available.

The CDFD bridge here is:
\begin{equation}
    C_{\mathrm{top}}(\mathbf{x},t)
    =
    C_{\mathrm{bulk}}+
    C_{\mathrm{boundary}}+
    C_{\mathrm{defect}}+
    C_{\mathrm{field}}.
\end{equation}
If topology protects a channel, the effective constraint along that channel is low while the surrounding bulk constraint remains high. If an external magnetic field writes chains in a ferrofluid, it changes $C_{\mathrm{field}}$ and therefore the macroscopic flow response.

The key falsification test is again incremental prediction. If a topological invariant, field geometry, or defect map already predicts the transport completely, CDFD is only a translation. If the adaptive constraint law predicts transition timing, hysteresis, or failure of protection, it earns a place.

\section{Additional Mystery Queue}

The following topics are worth keeping in the expandable spine. Some should remain short subsections until there is enough mathematics or data for a standalone paper.

\begin{center}
\begin{tabular}{L{0.22\textwidth}L{0.34\textwidth}L{0.30\textwidth}}
\toprule
\textbf{Mystery} & \textbf{CDFD bridge} & \textbf{Needed discriminant}\\
\midrule
High-$T_c$ superconductivity & phase stiffness and dissipative constraint decouple before full coherence & distinguish pairing scale from phase-locking scale\\
Strange metals & charge flux lacks ordinary quasiparticle constraint & transport scaling tied to adaptive scattering field\\
Quantum Hall and topological phases & transport protected by topological constraint channels & edge-state resilience predicted by constraint topology\\
Sonoluminescence & collapse focuses hydrodynamic flux through transient constraint shell & emission tied to maximum compression topology\\
Lightning and ball-lightning candidates & atmospheric charge flux constrained by plasma channel memory & persistence depends on channel constraint lifetime\\
Non-Newtonian fluids & imposed shear modifies the constraint field itself & shear-thickening/thinning predicted by $C(\dot{\gamma})$ dynamics, where $\dot{\gamma}$ is shear rate\\
Fracture and crack propagation & stress flux localizes along weakening constraint corridors & crack path follows high-$\nabla(C^{-1}\nabla\Phi)$ regions\\
Magnetorheological and ferrofluids & external field writes transport constraint topology & viscosity/yield stress follows field-imposed $C$ geometry\\
\bottomrule
\end{tabular}
\end{center}

\subsection{Priority Ranking}

Not all mysteries should be expanded equally. The best next candidates are the ones where data can be produced or found quickly:
\begin{enumerate}
    \item \textbf{Friction and fracture}: accessible laboratory systems with imaging and acoustic signals.
    \item \textbf{Plasma}: existing confinement and reconnection diagnostics can test drive-plus-topology predictors.
    \item \textbf{Superconductivity}: harder, but strategically powerful because CDFD must distinguish phase stiffness from dissipation.
    \item \textbf{Jamming and glasses}: excellent for memory and preparation-history tests.
    \item \textbf{Lightning and sonoluminescence}: interesting but should stay exploratory until the variables are better pinned down.
\end{enumerate}

\section{The Mujjabi Falsification Program}

Paper I introduced $M_s$ as a possible memory or recovery state of the vacuum
medium. That idea is not established by the numerical recovery of
$\chi=1/\alpha$, where un-subscripted $\alpha$ denotes the electromagnetic
fine-structure constant. It becomes scientific only if it creates discriminants that can
be compared against QED, general relativity, and established laboratory-domain
models. The tests below are therefore written as falsification targets, not as
claimed observations.

\subsection{Target parameters from the source notes}

The DOCX source notes singled out these experiments as high-priority tests. This
manuscript keeps descriptive labels, because the scientific content is the
measurable discriminant. The common parameter is the memory relaxation rate
\begin{equation}
    \lambda_M=\tau_M^{-1},
\end{equation}
where $\tau_M$ is the vacuum-memory recovery time proposed in Paper I. No
numerical value is assigned here. A value or bound must come from either a deeper
vacuum equation of state or a controlled null/result in the experiments below.

\begin{center}
\small
\begin{tabular}{L{0.25\textwidth}L{0.31\textwidth}L{0.30\textwidth}}
\toprule
\textbf{Test family} & \textbf{Primary target} & \textbf{Required control}\\
\midrule
Mujjabi Vacuum Hysteresis Test & history-dependent phase residual decaying with $\tau_M$ & nonlinear QED, plasma, gas, optics, thermal, and detector effects\\
Mujjabi Alpha-Jitter Test & residual alpha-sensitive transition shift scaling with flux intensity & standard QED, nuclear structure, thermal, magnetic, and systematic corrections\\
Mujjabi Rotating-Source Test & spin-correlated pressure-gradient signal & Newtonian gravity, vibration, magnetic coupling, electrostatics, air drag, and motor heating\\
\bottomrule
\end{tabular}
\end{center}

\subsection{Mujjabi Vacuum Hysteresis Test}

Let a high-field pulse produce a temporary overload functional
\begin{equation}
    \mathcal{H}(t)
    =
    \int_{\Omega}
    \max\!\left(0,\frac{|\bm{J}(\bm{x},t)|}{C(\bm{x},t)}-1\right)\,dV.
\end{equation}
If the memory kernel proposed in Paper I is physical, a second probe pulse sent
before the relaxation time $\tau_M$ has elapsed should see a residual response
that depends on the history of the first pulse:
\begin{equation}
    \Delta\phi_2
    =
    F(I_1,\Delta t,\tau_M)-F(0,\Delta t,\tau_M),
\end{equation}
where $I_1$ is the pump intensity and $\Delta t$ is the pump--probe delay.

\begin{prediction}[Vacuum hysteresis]
After all known nonlinear QED, plasma, gas, mirror, thermal, and detector
effects are controlled, a nonzero history-dependent phase residual should scale
with pump intensity and decay with a finite relaxation time if $M_s$ is a real
vacuum state variable.
\end{prediction}

\begin{criterion}[Vacuum-hysteresis falsification]
The vacuum-memory bridge fails if high-field pump--probe experiments find no
history-dependent residual beyond known QED and instrumental effects, or if any
residual does not scale with the overload functional $\mathcal{H}$.
\end{criterion}

\subsection{Quantum Decoherence as Vacuum Hysteresis}

The source notes propose a stronger bridge: decoherence may contain a
vacuum-memory component. This manuscript keeps that idea in testable form. In
standard open quantum systems, loss of interference is modeled by environmental
entanglement and dephasing channels. CDFD can add value only if, after those
channels are fitted, a residual depends on the prior loading history of the
experimental volume.

\paragraph{Observer-effect clarification.}
In this manuscript, observation means physical coupling, not conscious attention.
A detector, cavity wall, probe pulse, thermal bath, residual gas, vibration
mode, or readout circuit can change a quantum experiment because it interacts
with the system and leaves a record. A human observer may later inspect that
record, but the relevant physics is the coupling channel. The CDFD proposal below
is therefore limited to a possible history-dependent residual after ordinary
measurement, environmental, and instrumental channels have been modeled.

Let $V(t)$ denote an observed interference visibility, and let
$V_{\mathrm{std}}(t)$ be the visibility predicted by a calibrated standard
open-system model using the known environmental couplings. Define the residual
\begin{equation}
    R_{\mathrm{dec}}(t)
    =
    V(t)-V_{\mathrm{std}}(t).
\end{equation}
The CDFD claim is that a nonzero residual, if present, should scale with a
memory state $M_s-1$ generated by earlier flux loading:
\begin{equation}
    R_{\mathrm{dec}}(t,\Delta t)
    \propto
    \left[M_s(t)-1\right]e^{-\Delta t/\tau_M},
\end{equation}
where $\Delta t$ is the interval between the loading event and the coherence
measurement. This is not a replacement for standard decoherence theory. It is a
possible history-dependent correction that must vanish when the medium has
relaxed or when no overload history was applied.

\begin{prediction}[Decoherence-memory residual]
For otherwise identical quantum-interference experiments, a deliberately
pre-loaded high-flux history should change the subsequent visibility only if a
vacuum-memory term exists. The sign and magnitude of the residual must track
the loading intensity and decay with the same $\tau_M$ tested by the
pump--probe hysteresis protocol.
\end{prediction}

\begin{criterion}[Decoherence-memory falsification]
The decoherence-memory bridge fails if preparation-history changes produce no
reproducible residual after standard dephasing, thermal, vibration,
electromagnetic, detector, and gas effects are bounded, or if residuals do not
scale with the proposed memory variable.
\end{criterion}

\subsection{Mujjabi Alpha-Jitter Test}

The strongest publishable form of this idea is narrow. If the electromagnetic
fine-structure constant $\alpha=a/R$ is an
adaptive fixed point rather than a rigid input, extreme field
environments could in principle create a tiny effective displacement:
\begin{equation}
    \frac{\Delta\alpha}{\alpha}
    =
    \mathcal{F}\!\left(\frac{\Phi}{C},S,M_s\right),
    \qquad
    \mathcal{F}(0,1,1)=0.
\end{equation}
The natural search space is high-precision spectroscopy of highly charged ions,
strong-field QED systems, astrophysical spectra, or other environments where
standard calculations already predict the dominant shifts.

\begin{prediction}[Extreme-field alpha sensitivity]
If adaptive vacuum response is real, transition energies that depend strongly on
$\alpha$ should show a small residual correlated with controlled field intensity
or environmental flux after all standard QED, nuclear, thermal, and systematic
corrections are applied.
\end{prediction}

\begin{criterion}[Spectral falsification]
The bridge fails if high-precision spectra continue to match standard QED and
known systematics with no residual scaling with $\Phi/C$, or if any residual is
not reproducible across independent platforms.
\end{criterion}

\subsection{Mujjabi Rotating-Source Test}

The gravitational extension of CDFD is the weakest of the three vacuum-level
targets and should be treated accordingly. The hypothesis is that rapidly
organized angular momentum in dense matter may produce a secondary pressure or
constraint-gradient contribution:
\begin{equation}
    \Delta g_{\mathrm{CDFD}}
    \propto
    \nabla\!\left(C_{\mathrm{vac}}^{-1}\Phi_{\mathrm{rot}}\right).
\end{equation}
A defensible experiment would require shielded, balanced rotors, blind reversal
of spin direction, vibration isolation, magnetic/electrostatic controls, and an
atom-interferometric or torsion-balance readout.

\begin{prediction}[Rotating-source residual]
After conventional Newtonian, relativistic, vibration, thermal, electromagnetic,
and air-coupling effects are subtracted, any remaining spin-correlated signal
should scale with angular momentum and source geometry, not merely with motor
power or mechanical noise.
\end{prediction}

\begin{criterion}[Rotating-source falsification]
The rotating-mass bridge fails if controlled experiments find no spin-correlated
gravitational residual, or if observed residuals track known engineering
systematics rather than the predicted constraint-gradient geometry.
\end{criterion}

\subsection{Experimental priority}

The priority order should be conservative:
\begin{enumerate}
    \item high-frequency vacuum hysteresis, because it directly targets $M_s$;
    \item extreme-field spectroscopy, because it tests whether $\alpha$ behaves
          as a rigid constant or an adaptive fixed point;
    \item rotating-source gravity, because the expected signal is most vulnerable
          to mundane systematics and therefore needs the strictest controls.
\end{enumerate}
None of these tests is evidence by itself. Their value is that they make the
non-equilibrium extension of Paper I falsifiable.

\section{Unified Falsification Program}

\begin{center}
\small
\begin{tabular}{L{0.22\textwidth}L{0.30\textwidth}L{0.34\textwidth}}
\toprule
\textbf{Domain} & \textbf{Primary measurement} & \textbf{Failure condition}\\
\midrule
Friction & contact-resistance noise, acoustic emission, real contact area & no residual memory/geometric predictor after rate-and-state modeling\\
Superconductivity & phase stiffness, resistance, vortex motion, disorder & CDFD variables do not improve transition or vortex predictions\\
Plasma & current profile, magnetic topology, turbulence, confinement time & disruptions depend only on known scalar parameters\\
Turbulence & local Reynolds fields, boundary-layer stress, coherent structures & no local constraint field improves transition localization\\
Jamming/glass & packing, contact network, relaxation time, avalanche sizes & preparation history adds no predictive information\\
\bottomrule
\end{tabular}
\end{center}

The correct standard is incremental predictive value. CDFD does not win by renaming known physics. It wins only if the $\Phi/C$ mapping and adaptive-constraint law predict where a transition occurs, how it scales, or why two apparently similar systems diverge.

\section{Reproducibility Artifacts}

Paper XII now has its own output set:
\begin{center}
\small
\begin{tabular}{L{0.42\textwidth}L{0.42\textwidth}}
\toprule
\textbf{Artifact} & \textbf{Purpose}\\
\midrule
\url{fig12_0_transport_mystery_map.pdf} & visual flux/constraint map for the included mysteries\\
\url{fig12_1_friction_stick_slip.pdf} & adaptive-memory stick-slip toy model\\
\url{fig12_2_superconductivity_reallocation.pdf} & dissipative vs phase-constraint transition model\\
\url{fig12_3_plasma_rupture_index.pdf} & drive/topology rupture-index field\\
\url{fig12_4_jamming_memory.pdf} & preparation-memory avalanche distribution\\
\url{checks_summary_XII_gate.csv} & pass/fail summary for generated outputs\\
\bottomrule
\end{tabular}
\end{center}

The generator is intentionally modest. The outputs are not evidence for the physical theory; they are a precise statement of what the proposed theory means operationally. This matters because a falsifiable wrong model is more useful than an unfalsifiable analogy.

\section{Relation to Papers I--XI}

Paper I introduced the physical-vacuum intuition using vortex stability. Papers II--X developed the mass-spectrum and vacuum-EOS ladder. Paper XI used that ladder to frame cosmological hypotheses. Paper XII is different: it is a proposed laboratory bridge. It asks whether the same flux-constraint vocabulary can organize measurable collective states without leaning on cosmology.

The strongest reason to keep Paper XII is strategic. A reader may reject the dark-sector bridge, but still test frictional memory, vortex pinning, plasma rupture, or jamming as constrained transport. These are nearer to experiment, easier to falsify, and useful for sharpening the whole theory.

\section{Conclusion}

Friction, superconductivity, plasma confinement, turbulence, and jamming are not random additions to the physics series. They are the everyday places where transport mystery is most visible. A surface locks then slips. A material loses resistance but gains phase rigidity. A plasma confines then reconnects. A fluid flows smoothly then becomes turbulent. A dense material behaves solid until it avalanches. In CDFD language, these are all candidates for adaptive constraint transitions.

The claim remains conditional. If the measurable flux and constraint fields do not improve prediction beyond established models, Paper XII is only a metaphor. If they do, these laboratory mysteries become the bridge between the formal CDFD vacuum papers and practical physics.

\section*{Supplementary Material}

This manuscript is accompanied by release-archive supplementary material in
the canonical Part I tree. The relevant paper-local Python scripts, numerical checks, generated figure outputs,
tabulated data, figure assets, and environment notes are maintained under
\texttt{notebooks/}, \texttt{outputs/}, \texttt{figures/}, and
\texttt{REPRODUCIBILITY.md} in \texttt{Part\_I\_Fundamental\_Physics/}.


\section*{Data and Code Availability}

The release-archive supplementary material is included in the archived Part I
release on Zenodo at \url{https://doi.org/10.5281/zenodo.20250820}. The current version DOI is assigned by Zenodo after each GitHub release. Zenodo is the
permanent release archive, and the public source archive is available on GitHub at
\url{https://github.com/bampita-bico/CDFD-Part-I-Release}. The broader public CDFD
Runtime is separate reusable software infrastructure; the numerical values reported
here are reproduced from this paper's Supplementary Material.

\section*{Acknowledgments}

The author used AI-assisted tools for formatting checks and code-verification
support. The scientific claims, interpretations, and final manuscript decisions
are the author's responsibility.

\section*{Declarations}

\subsection*{Funding}
This research received no external institutional funding.

\subsection*{Author Contributions}
Steve Bico Mujjabi: conceptualization, methodology, formal analysis, software,
validation, visualization, writing--original draft, and writing--review and
editing.

\subsection*{Conflicts of Interest}
The author declares no conflicts of interest.

\section*{References}
\bibliographystyle{apsrev4-2}
\bibliography{references}

\end{document}
